\section{Conclusão}

O projeto confirmou a eficácia de sistemas biométricos distribuídos para controle de acesso, integrando processamento local e nuvem para garantir segurança e performance. A arquitetura centralizada assegura uma gestão auditável e escalável, atendendo aos requisitos de alta confiabilidade corporativa.

No âmbito da visão computacional, o uso do \textit{ArcFace} com processamento local reduziu gargalos de latência e tráfego de rede. Embora o sistema tenha superado as limitações de soluções que dependem de imagens brutas, observou-se a ocorrência de falsos negativos. Tais erros foram atribuídos a variações nos vetores de características, causadas por diferenças na iluminação e na qualidade do sensor da câmera entre a captura em tempo real e o armazenamento no banco de dados.

Para trabalhos futuros, recomenda-se a implementação de detecção de vivacidade (\textit{liveness detection}) para elevar a segurança contra fraudes. Adicionalmente, a otimização dos modelos para hardware de borda e o refinamento dos critérios de qualidade da imagem permitirão escalar a solução para mais dispositivos, mitigando as variações que impactam a precisão do reconhecimento.-